% =====================================================================
%  Kapitel 4 – Verfahren und Methoden (~6–7 Seiten)
% =====================================================================
\chapter{Verfahren und Methoden}
\label{kap:methoden}

\section{Systemarchitektur}
\label{sec:architektur}

Abbildung~\ref{fig:architektur} zeigt den Aufbau des Systems.

\begin{figure}[htbp]
  \centering
  \begin{tikzpicture}[
      node distance = 1.2cm and 1.4cm,
      block/.style = {draw, rounded corners, align=center,
                      minimum width=2.8cm, minimum height=1cm, font=\small},
      >={Stealth},
    ]
    \node[block] (sens)  {Sensoren\\{\footnotesize Thermal, Velostat}};
    \node[block, right=of sens] (board) {Neffy Board\\{\footnotesize ESP32-S3}};
    \node[block, right=of board] (pi)   {Raspberry Pi};
    \node[block, above right=0.4cm and 1.4cm of pi] (motor) {Motor};
    \node[block, below right=0.4cm and 1.4cm of pi] (disp)  {Display};
    \node[block, below=of board] (led)  {\ac{LED}s\\{\footnotesize WS2812B}};

    \draw[->]  (sens)  -- (board);
    \draw[<->] (board) -- (pi);
    \draw[->]  (pi)    -- (motor);
    \draw[->]  (pi)    -- (disp);
    \draw[->]  (board) -- (led);
  \end{tikzpicture}
  \caption{Systemarchitektur von Neffy~3.0 (Entwurf)}
  \label{fig:architektur}
\end{figure}
\todo{Architektur prüfen: Wer steuert den Motor? Schnittstelle ESP32 -- Pi (UART/USB/...)?}

\section{Hardware und Sensorik}
\label{sec:hardware}

\subsection{Anwesenheitserkennung}
\label{sec:hw_anwesenheit}

\todoi{Omron D6T-1A-02: Funktionsprinzip, Anschluss über \ac{I2C},
Datenblatt zitieren. Final integriert?}

\subsection{Auswahl des Drucksensors}
\label{sec:hw_drucksensor}

\todoi{Muster: Problem $\rightarrow$ Optionen $\rightarrow$ Kriterien
$\rightarrow$ Entscheidung $\rightarrow$ Begründung. Getestete Varianten
(u.\,a. \ac{FSR}, \dots), Beobachtungen, Entscheidung für Velostat.}

\begin{table}[htbp]
  \centering
  \caption{Bewertung der getesteten Drucksensoren}
  \label{tab:bewertung_drucksensor}
  \small
  \begin{tabular}{@{}p{3.4cm} p{3.2cm} p{3.2cm} p{3.2cm}@{}}
    \toprule
    Kriterium & \ac{FSR} & Velostat & \todoi{weitere Variante} \\
    \midrule
    \todoi{Kriterium 1} & & & \\
    \todoi{Kriterium 2} & & & \\
    \bottomrule
  \end{tabular}
\end{table}

\subsection{Messaufbau und Schaltung}
\label{sec:hw_schaltung}

\todoi{Schaltung Velostat am ESP32-S3 (z.\,B. Spannungsteiler, Widerstandswert,
ADC-Pin), Messaufbau, Foto.}

\section{Konstruktion des Gehäuses}
\label{sec:gehaeuse}

\todoi{Eiförmige Schale, gespiegelte Hälften, Teilung in druckbare Teile mit
abnehmbarem Deckel; \ac{CAD} in SolidWorks.}

\begin{figure}[htbp]
  \centering
  \begin{subfigure}[b]{0.31\linewidth}
    \centering
    \platzhalter[\linewidth]{CAD-Modell}
    \caption{CAD-Modell}
  \end{subfigure}\hfill
  \begin{subfigure}[b]{0.31\linewidth}
    \centering
    \platzhalter[\linewidth]{Teilung}
    \caption{Teilung}
  \end{subfigure}\hfill
  \begin{subfigure}[b]{0.31\linewidth}
    \centering
    \platzhalter[\linewidth]{Druckteil}
    \caption{Druck}
  \end{subfigure}
  \caption{Entwicklungsschritte des Gehäuses}
  \label{fig:gehaeuse_schritte}
\end{figure}

\subsection{Teilung und Ausrichtung}
\label{sec:gehaeuse_teilung}

\todoi{Passstifte Ø5/Ø6\,mm, Rundloch-plus-Langloch-Prinzip -- Begründung.}

\subsection{Toleranzen im \ac{FDM}-Druck}
\label{sec:gehaeuse_toleranz}

\todoi{Testkörper zur empirischen Ermittlung der Passungen: Maße, Ergebnisse.}

\subsection{Fertigung und Nachbearbeitung}
\label{sec:gehaeuse_fertigung}

\todoi{Drucker, Material, Druckzeiten, Schleifen, Lackieren.}

\section{Software}
\label{sec:software}

\subsection{Signalverarbeitung und Atemerkennung}
\label{sec:sw_atem}

\todoi{Flankenerkennung (aktive Variante), 20-s-Fenster, mittlerer Atemrhythmus.}

\subsection{Softwarestruktur und Ablage}
\label{sec:sw_struktur}

\todoi{Modulstruktur (.cpp/.h), PlatformIO, FastLED; Repository-Link, Build-Anleitung.}

\subsection{Einsatz von KI-Werkzeugen}
\label{sec:sw_ki}

\todoi{Transparent benennen, wofür KI-Werkzeuge genutzt wurden.}
